% Folder 83, batch 05 (archivists' pages 81-100).
% Pass: Opus 5.5 (claude-opus-5-5), 2026-10-10 — first pass, unchecked against the pages by a human.
%
% Page 93 is a typescript leaf of his non-mathematical writing (« PU 65 »,
% on abstraction and concepts), cancelled and annotated in his hand; not the
% mathematics in hand, skipped. Page 100 is a blank sheet; skipped.
% Page 92's verso carries a typescript; only his handwritten mathematics on
% it (multiplication table and trace matrix) is transcribed.
% Pages 81-99: the algebra M(b,c;b',c';t) generated by u, v with Tr u = b,
% det u = c, Tr v = b', det v = c', Tr(uv) = t; Hamilton-Cayley; the
% anharmonic invariant rho; Q = t^2 - bb't + b^2c' + b'^2c - 4cc'; when is
% M Azumaya (theorem p. 98: iff det restricted to k+ku+kv non-degenerate).
\documentclass[11pt,a4paper]{article}
\input{../preamble/grothendieck}

\folder{83}
\batch{5}
\pages{81}{100}
\foldertitle{[Icosaèdres, bi-icosaèdres et algèbres d'Azumaya de rang 4] : notes manuscrites (1982-1983, 1986, s.d.).}
\dating{1982-1986}
\watermark{Édition de démonstration}

\begin{document}

\page{81}
\note{L'argument vient de la page précédente, hors du lot : les relations (1), (2), (3) (Hamilton--Cayley, « H.C. ») et la numérotation des formules y ont été posées.}
de la première des relations
\[
(5)\qquad \begin{cases} \operatorname{Tr} 1 = 2 \\ \det 1 = 1 \end{cases}
\]
\add{Dont la 2°} \struck{lesquelles} se déduit \struck{\ill{}} de (2) en y faisant $u=1$ (et $\det 1 = 1$ est donné \uncertain{comme connu}).

En polarisant (2), on trouve pour $u, v \in A$ (tenant compte de (3))
\[
(6)\qquad uv + vu - (\operatorname{Tr} u)v - \operatorname{Tr}(v)u + \bigl(\operatorname{Tr} u \operatorname{Tr} v - \operatorname{Tr}(uv)\bigr)1 = 0
\]
soit
\[
(6')\qquad uv + vu = \operatorname{Tr}(v)u + (\operatorname{Tr} u)v + \bigl(\operatorname{Tr}(uv) - \operatorname{Tr} u \operatorname{Tr} v\bigr)1 .
\]
\note{la page porte, en travers de (6) et (6'), plusieurs lignes obliques serrées que la lecture ne démêle pas ; on y lit « (6') résulte de \ill{} », « NB Cette relation \ill{} », « \uncertain{identité} », « (2), (3), (6) », et un renvoi « (6'') ». La ligne biffée qui suit (6') commence par \struck{On peut \uncertain{déduire}}.}
\marginal{\ill{} $(6'')$ \ill{}}

$k.1 + ku + kv + \ill{}$ \ill{} \struck{\uncertain{opposé}} opposé \ill{} $kl + ku + kv$.

\textbf{Proposition}. \struck{Soit} Soit $A$ \ill{} (1) (2) (3) (donc (6')). Alors $\forall u, v \in A$, \struck{$k + ku + kv + k.uv$} le sous-module engendré par $1, u, v, uv$ est une sous-algèbre.

Il suffit de voir que \struck{le produit de}

\page{82}
deux \struck{éléments} parmi les éléments $1, u, v, uv$ est dans ce sous-module. Il suffit de le voir pour deux tels éléments distincts de $1$. \struck{On \ill{}} \note{suit un tableau biffé : $u.u = \det(u)1 + (\operatorname{Tr} u)u + \ldots$, $u.v = 0.1 + 0.u + 0.v + 1uv$, $u.uv = \ldots$}

\[
u^2,\ v^2,\ uvuv,\ uv,\ u^2v,\ vu,\ vuv,\ uvu,\ uv^2
\]
Or
\[
\begin{aligned}
u^2 &= \operatorname{Tr}(u).u - \det(u).1 && \text{(2) pour } u\\
v^2 &= \operatorname{Tr}(v).v - \det(v).1 && \text{(2) pour } v\\
uv &= 1.uv && \text{tautol.}\\
vu &= -1.uv + \operatorname{Tr}(v)u + \operatorname{Tr}(u)v + \bigl(\operatorname{Tr}(uv) - \operatorname{Tr} u \operatorname{Tr} v\bigr)1 && \text{(relation (6'))}
\end{aligned}
\]

Ce sous-module $V$ est stable par multiplication \add{à gauche} par $1, u, v$, \struck{$uv$ \ill{}} \struck{\ill{} \ill{} $uv$, et} (\uncertain{compte} 1). Donc il suffit de voir que $ux, vx \in V$, si $x \in V$. Il suffit de prendre pour $x$ un des générateurs $1, u, v, uv$. Le cas $1$ étant trivial, il reste
\[
u^2,\ uv,\ u^2v,\ vu,\ v^2,\ vuv = \struck{\ill{}}\,(vu)v .
\]
Pour $u^2, v^2$ on utilise H.C. (3), pour $uv$ c'est tautologique, pour $u^2v$ cela résulte de l'expression de $u^2$ (H.C.), pour $vu$ on utilise (6'), qui s'écrit
\[
\text{\struck{$vu = \ill{} + \operatorname{Tr}(v).u +$}}
\]
\[
(6'')\qquad vu = \bigl(\operatorname{Tr}(uv) - \operatorname{Tr} u \operatorname{Tr} v\bigr)1 + \operatorname{Tr}(v)u + \operatorname{Tr}(u)v - uv .
\]
Enfin pour \struck{\ill{}} $vuv = (vu)v$, on multiplie la relation (6'') à droite par $v$, on utilise H.C. pour $v$ (pour se débarrasser des termes en $v^2$).

\page{83}
Finalement, on trouve une « table de multiplication » pour une $k$-algèbre libre de rang 4 (dont $V$ sera quotient) en termes de 5 éléments de $k$
\[
(7)\qquad \begin{cases} b = \operatorname{Tr} u,\ c = \det u \\ b' = \operatorname{Tr} v,\ c' = \det v \\ t = \operatorname{Tr}(uv) \end{cases}
\]
(permet de décrire $k[u] = k[T]/(T^2 - bT + c)$ (si $k \cap k.u = \{0\}$)) (id pour $k[v]$)

avec comme base $1, u, v, uv$ (les \uncertain{analogues} qu'on considère \uncertain{étant} de la \uncertain{nouvelle} algèbre). On appelle cette algèbre
\[
(8)\qquad M(b, c;\, b', c';\, t) .
\]
On vérifie qu'elle existe (i.e. indépendamment \uncertain{de la donnée} du (7)) pour toute la donnée \struck{\ill{}}
\[
(9)\qquad (b, c, b', c', t) \in k^5 .
\]
En fait, on a une \uncertain{alg.} universelle sur
\[
\mathbf{Z}[B, C, B', C', T]
\]
où $B, C, B', C', T$ sont des indéterminées.

\emph{Vérification}. On écrit une table de multiplication explicite en utilisant H.C. pour $u$ et $v$ et (6) alias (6') (pour $(u, v)$),

\page{84}
et on vérifie l'associativité. De plus, on prolonge $\operatorname{Tr}$ à $V$ par linéarité, et $\det$, \struck{pour} d'abord sur $\underbrace{k[u]}_{A_1} \subset M$ et $\underbrace{k[v]}_{A_2} \subset M$ (où c'est la norme \struck{$\det(u)'$} des algèbres quadratiques, $\operatorname{Tr}$ y étant la trace) puis sur
\[
M = k[u] \oplus k[v]
\]
\note{au-dessus et à côté de cette égalité : « iso » et « (\uncertain{à cause des} \ill{} \uncertain{identité}) », qu'on ne démêle pas.}
par
\[
(10)\qquad \det(u_1 + v_1) = \underbrace{\det u_1}_{N_{A_1/k}(u_1)} + \underbrace{\det v_1}_{N_{A_2/k}(v_1)} + \operatorname{Tr} u_1 \operatorname{Tr} v_1 - \operatorname{Tr}(u_1v_1)
\]
\note{le premier terme est écrit $\det(x_1 + v_1)$ sur la page, avec une surcharge ; on lit $u_1$ par la suite de la formule.}

Ceci dit, on « vérifie » \struck{\ill{}} que $\det$ est quadratique et multiplicative, que l'on a la relation de Hamilton--Cayley et les relations (3).

En fait, je n'ai pas eu à faire la vérification, car \struck{je sais} il suffit de le faire (par principe de prolongement des id. alg.) \struck{pour un} sur un \ill{} \ill{} \uncertain{dense}, \ill{} \ill{}
\[
\delta \overset{\text{déf}}{=} b^2 - 4c \text{ inv.}, \qquad \delta' \overset{\text{déf}}{=} b'^2 - 4c' \text{ inv.}
\]
(i.e. $A_1, A_2$ étales sur $k$), et on voit \add{pour $b, c, b', c'$ donnés, ainsi,} que \struck{l'on peut réaliser (pour} \struck{alors (pour \ill{}} \struck{la répartition $(b, c$}

\page{85}
la répartition $(b, c, b', c', t)$ « pour presque tout $t$ » \add{et toute $c$} \add{par deux éléments $u, v$} donnant l'alg. de Matrices \add{$M = M_2(k)$}. En effet, \struck{pour \ill{}} \struck{\ill{} $u, v$} on localise pour que $A_1, A_2$ soient diagonalisables ($\simeq k \times k$), on réalise comme sous-algèbres de \ill{} \ill{}, \struck{\ill{} \ill{} de \ill{}} pour plonger les $A$, i.e. par $u \in A$, $v \in A$.

Bien sûr, on n'a \struck{\ill{} pu \ill{}} pu préjuger de la valeur de $\operatorname{Tr} uv$. Mais \struck{considérons, pour}

\emph{Lemme}. Soient $u, v \in A$ ($\simeq M_2(k)$) qui \struck{\ill{}} engendrent $k[u]$, $k[v]$ étales de rang 2 sur $\operatorname{Spec} k$. Considérons pour $\Theta$ variable dans $A^*$
\[
\operatorname{Tr}(u\Theta(v)), \qquad \Theta(v) \overset{\text{déf}}{=} \Theta v \Theta^{-1} .
\]
On trouve un morphisme de schémas
\[
\check{\mathbb{V}}(A) \longrightarrow \mathbb{E}^1_k \;(= \operatorname{Spec} k[T])
\]
\emph{dominant}.
\marginal{NB $\Theta \mapsto \Theta(v)$, \ill{} \ill{} \ill{} $u, v$ \ill{}}

En fait, on trouve une fonction rationnelle en $b, c, b', c', t$, \ill{} \ill{} discriminant $\ldots t^2 - 4cc'$ ($= \operatorname{Tr}(uv)^2 - 4\det uv = \delta(uv)$, discriminant de $uv$)

\page{86}
(\uncertain{invariant} \emph{anharmonique} de $k[u]$ et $k[v]$ dans \uncertain{alg.} d'Azumaya) de
\[
\rho = f(b, c, b', c';\, t) = \frac{P(b, c, b', c', t)}{Q(\ldots)}
\]
telles que la condition nécessaire et suffisante pour que \struck{\ill{}} $b, c, b', c', t$ \add{se réalisent} \struck{soit dans \ill{} d'Azumaya}, \struck{\ill{} que} (en \uncertain{supposant} $\delta\delta'$ inversibles), c'est que $\rho \neq 2$ \add{partout}, i.e.
\[
P - 2Q \neq 2 \qquad \text{\uncertain{partout}}
\]
\note{ainsi sur la page ; on attend $P - 2Q$ partout inversible.}

En fait, quand ils sont réalisés dans \uncertain{alg.} d'Azumaya, \struck{ils définissent} ils définissent \struck{deux des} chacun un diviseur de degré 2, \struck{\ill{}} \uncertain{dans} les schémas projectifs associés à $A$, et si leurs supports sont disjoints, il y a un invariant anharmonique $\lambda$ qui n'est défini que modulo l'involution $\lambda \mapsto \lambda^{-1}$ (\uncertain{défini} \ill{})
\[
\Bigl(\rho = \lambda + \text{\struck{$\ill{}$}}\,\lambda^{-1} \qquad \text{d'où} \quad \lambda^2 - \rho\lambda + 1 = 0
\]
(et comme $\lambda \neq 0, 1$, on dira \ill{} $\rho \neq 2$ \ldots

\page{87}
J'ai fait le calcul de $\rho$ en termes de $b, c, b', c', t$) dans le cas où $b = \operatorname{Tr} u$ et $b' = \operatorname{Tr} v$ sont nuls, \struck{on trouve} (supposant que $u, v$ plongés dans $A$, et \uncertain{Alg.} d'Az.). On trouve donc
\[
\rho = 2\,\frac{t^2 + 4cc'}{t^2 - 4cc'}
\qquad\Bigl|\qquad
\begin{gathered}\text{NB } t^2 - 4cc' = \delta(uv)\\ \text{si } b = b' = 0, \text{ discriminant de } uv\end{gathered}
\]
d'où
\[
\frac{\rho}{2} - 1 = 8\,\frac{cc'}{t^2 - 4cc'} .
\]
On en déduit le cas général [Il suffit de le faire sur $\mathbf{Z}[\tfrac12][B, C, B', C', T]$]. On trouve
\[
(11)\qquad \rho = 2\,\frac{(2t - bb')^2 + \delta(u)\delta(v)}{(2t - bb')^2 - \delta(u)\delta(v)}
\]
\[
= \frac{2t^2 - 2bb't + \bigl(b^2b'^2 - 2(b^2c' + b'^2c) + 8cc'\bigr)}{t^2 - bb't + (b^2c' + b'^2c - 4cc')}
\]
\note{sous $\rho$, une flèche vers « $\delta u$ si 2 inv. dans $k$ » ; sous le signe $=$, « \uncertain{$\delta h$} \ill{} \ill{} ».}

[En \uncertain{car.} 2, cela donne
\[
\rho = \frac{(bb')^2}{t^2 - bb't + (b^2c' + b'^2c)} \; .]
\]

\[
(**)\qquad \rho_0 = \rho - 2 = \frac{b^2b'^2 - 4(b^2c' + b'^2c) + 16cc'}{t^2 - bb't + (b^2c' + b'^2c - 4cc')} = \frac{A(b, b', c, c')}{B(t; b, b', c, c')}
\]
\marginal{invariant anh. rectifié (pour $\rho_0$)}

Donc
\[
\rho\, B(t; b, b', c, c') - A(b, b', c, c') = 0
\]
\note{ainsi sur la page ($\rho$ et non $\rho_0$) ; le numéro $(**)$ surcharge un autre numéro.}

\page{88}
$A_1$ et $A_2$ \add{diagonales} engendrant $M_2$. Cela se voit au moins loc. \uncertain{supposé} (en prenant à la \struck{solution} rac. des solutions de
\[
\lambda^2 - \rho\lambda + 1 = 0 \qquad \text{\struck{\ill{}}}
\]
d'abord \uncertain{posé}, à cause de $b^2 - 4c$ \struck{\ill{}},
on peut trouver loc. \uncertain{supposé} (prenant au rac. des solutions de $x^2 - bx + c = 0$),
$u$ dans $A_1$ correspondant :
\[
\operatorname{Tr} u = b, \quad \det u = c
\]
et item $v$ dans $A_2$
\[
\operatorname{Tr} v = b', \quad \det v = c' .
\]
Soit \struck{\ill{}}
\[
\tau = \operatorname{Tr} uv
\]
alors $\tau$ satisfait la relation analogue à (11') avec $t$ remplacé par $\tau$.

Or on trouve qu'\add{une} équation en $t$
\[
t^2 - bb't + \Bigl(b^2c' + b'^2c - 4cc' - \frac{\delta\delta'}{\rho_0}\Bigr)
\]
(où $\delta = b^2 - 4c$, $\delta' = b'^2 - 4c'$), il y a \struck{\ill{}} \struck{juste} (sur un corps) \struck{au plus} 2 solutions, \struck{disons} $t, t'$ (\uncertain{éventuellement} égales) satisfaisant
\[
t + t' = bb' \qquad (\tau = t') .
\]
Si $t = \tau$, \uncertain{on a} \uncertain{gagné}. Sinon, on remplace $u$ par

\page{89}
\note{cette page ne continue pas la phrase de la page 88 (« on remplace $u$ par \ldots »), dont la suite est en haut de la page 90 ; elle reprend (11') et énonce la proposition qui en sort.}
\[
(11')\qquad
\begin{cases}
\rho_0 = \rho - 2 = \dfrac{\delta\,\delta'}{t^2 - bb't + (\delta c' + \delta'c + 4cc')}\\[2ex]
\phantom{\rho_0 = \rho - 2} = \dfrac{(b^2 - 4c)(b'^2 - 4c')}{\underbrace{t^2 - bb't + b^2c' + b'^2c - 4cc'}_{Q\,=\,Q(t;\, b, b', c, c')}}
\end{cases}
\]
où $\delta = \delta(u) = b^2 - 4c$, $\delta' = \delta(v) = b'^2 - 4c'$.
\note{au numérateur de la première ligne, deux arguments de $\delta$ et $\delta'$ sont noircis ; le $4cc'$ du dénominateur surcharge un autre symbole.}

On trouve donc :

\textbf{Proposition}. Soient $k$ un anneau, et $b, c, b', c', t \in k$, tels que
\[
(b^2 - 4c)(b'^2 - 4c')(t^2 - bb't + b^2c' + b'^2c - 4cc')
\]
soit inv. dans $k$. Alors $M[b, c, b', c', t]$ est une alg. d'Azumaya \add{\uncertain{déf.}} (et $k[\underline{u}]$, $k[\underline{v}]$ en sont des sous-algèbres quadratiques étales, i.e. engendrant la \uncertain{conclusion} « en position générale »), \struck{\ill{}} dont les éléments \ill{} \uncertain{sont} \ill{}, \uncertain{algèbre} associative et satisfaisant les relations (2) (3) pour $\operatorname{Tr}$, $\det$ définis comme plus haut\ldots

\marginal{OPS $k$ corps \ill{}, \ill{} \ill{} des alg. \ill{} \uncertain{étendons} \ill{} \ill{}}
\emph{Dém}. On regarde $\rho_0$ comme défini plus \ill{} \uncertain{par} (11'). Comme $\rho_0$ \struck{$\neq 0$} \ill{} i.e. $\rho$ partout $\neq 2$, cela correspond à une position relative \ill{} \ill{} $\mathrm{SL}(2)$ (ou du \ill{} \ill{}) des tores des \ill{} \uncertain{alg.} quadratiques étales $A_1, A_2$ dans \ill{} \ill{} $M_2$) \ill{} \emph{position générale}. --- Cela signifie aussi que

\page{90}
\note{suite de la fin de la page 88.}
\[
u' = (\operatorname{Tr} u)1 - u = b.1 - u
\qquad
\begin{cases} \operatorname{Tr} u' = \operatorname{Tr} u \\ \det u' = \det u \end{cases}
\]
alors
\[
t' = \operatorname{Tr} u'v = \operatorname{Tr}\bigl((b - u)v\bigr) = b\operatorname{Tr}(v) - \operatorname{Tr}(uv) = bb' - \tau
\]
donc \struck{\ill{}} $t' = t$, on gagne !

Que se passe-t-il si
\[
(12_1)\qquad (b^2 - 4c)(b'^2 - 4c') \in k^*
\]
\struck{mais} \add{et} si $M(b, c, b', c';\, t)$ est azumayenne, je dis que
\[
(12_2)\qquad \underbrace{t^2 - bb't + b^2c' + b'^2c - 4cc'}_{Q(t;\, b, c, b', c')} \in k^* \;?
\]
\struck{Je dis que dans $M(b, c, b', c';\, t)$ \ill{}} \struck{\ill{} azumayenne}. En effet, \struck{s'il était}
comme $A_1 = k[\underline{u}]$ et $A_2 = k[\underline{v}]$ \struck{\ill{}} sont étales, et en position générale, il y aurait par la théorie \struck{la formule} \uncertain{générale} un invariant $\rho_0$, relié à $b, c, b', c', t$ par (11'), \uncertain{i.e.}
\[
\text{\struck{$\rho_0(t^2 - bb't$}} \quad \rho_0\, Q = \delta\delta'
\]
\uncertain{comme} $\delta\delta' \in k^*$, on aura $Q \in k^*$.

\textbf{Remarque}. Le discriminant de $Q$ (en tant que polynôme en $t$) est

\page{91}
\[
\Delta(Q) = (bb')^2 - 4(b^2c' + b'^2c - 4cc')
\]
On trouve
\[
(13)\qquad \Delta(Q) = (b^2 - 4c)(b'^2 - 4c') = \delta\delta'
\]
donc en fait, si \struck{\ill{}} comme \ill{} \ill{} $\delta\delta' \in k^*$, donc $\Delta(Q) \in k^*$ (supposé), i.e. « racines distinctes ». Cela signifie aussi que
\[
\text{\struck{$t \neq bb' - t$ \ i.e. \ $2t \neq bb'$}}
\]
\struck{$2\operatorname{Tr}(uv) =$}
\[
(14)\qquad \delta\delta' \in k^* \Longrightarrow 2t \neq bb' \ \text{\struck{\uncertain{pour tout} $t$}} \qquad \text{i.e. } 2t - bb' \in k^*
\]
\note{la ligne biffée et l'implication sont reliées par des traits ; on lit l'implication comme son remplacement.}

\struck{\ill{}} Le rev. des solutions de
\[
Q(t;\, b, c, b', c') = 0
\]
(\uncertain{en} $t$, pour $b, c, b', c'$ fixés) est étale (quadratique) sur $S = \operatorname{Spec} k$.

\textbf{Cor.} (?) Soient $u, v$ dans \add{une} alg. d'Azumaya $A$ \add{(de rang 4)}, tels que $\delta(u), \delta(v) \in k^*$, \struck{\ill{}} et $u, v$ partout \uncertain{linéairement indép.} mod $1$. \struck{tels que} Si $Q(u, v) = Q(t;\, b, c, b', c') = 0$, \struck{alors} il \ill{} \ill{} \ill{} \ill{}.
\marginal{\uncertain{Absurde}, \ill{} \ill{} \uncertain{donc} \ill{} \ill{} \ill{}}
\note{le corollaire, la marge oblique et la ligne biffée qui suit sont en partie barrés.}

\struck{$u, v$ \uncertain{appartiennent}} Tels que $u, v$ appartiennent à une
\struck{\ill{}} sous-alg. de Borel (qui n'est \uncertain{autre} alors que $k \oplus ku \oplus kv$) \struck{\ill{}} $uv \in k.1 + ku + kv$ \ldots

\page{92}
\[
\begin{array}{c|c|cccc}
 & 1 & \underline{u} & \underline{v} & \underline{uv}\\ \hline
\underline{1} & \underline{1} & \underline{u} & \underline{v} & \underline{uv}\\
\underline{u} & \underline{u} & -c + b\underline{u} & \underline{uv} & -c\underline{v} + b\underline{uv}\\
\underline{v} & \underline{v} & (t - bb')\underline{1} + b'\underline{u} + b\underline{v} - \underline{uv} & -c'\underline{1} + b'\underline{v} & -bc'\underline{1} + c'\underline{u} + t\underline{v}\\
\underline{uv} & \underline{uv} & -b'c\underline{1} + t\underline{u} + c\underline{v} & -c'\underline{u} + b'\underline{uv} & t\underline{uv} - cc'
\end{array}
\]
\note{table de multiplication de $M(b, c;\, b', c';\, t)$ dans la base $1, u, v, uv$ (facteur de gauche en ligne).}

En prenant les traces, on a la matrice
\note{une première matrice, à moitié remplie et en partie noircie, est barrée de traits obliques ; on en lit la première colonne $2, b, b', t$, la première ligne $2, b, b', t$, et $-2c + b^2$, $-2c' + b'^2$, $-b'c + bt$, $-bc' + b't$, $t^2 - cc'$. Elle est reprise proprement ci-dessous.}
\[
\begin{pmatrix}
2 & b & b' & t\\
b & b^2 - 2c & t & bt - b'c\\
b' & t & b'^2 - 2c' & b't - bc'\\
t & bt - b'c & b't - bc' & t^2 - cc'
\end{pmatrix}
\]
\marginal{dét. de $\alpha, \beta \mapsto \operatorname{Tr}(\alpha\beta)$ (par rapport à la base $1, u, v, uv$)}
\note{l'élément en bas à droite est écrit $t^2 - cc'$ ; on attendrait $\operatorname{Tr}((uv)^2) = t^2 - 2cc'$.}
\[
\begin{pmatrix}
4 & 2b & 2b'\\
2b & b^2 & bb'\\
2b' & bb' & b'^2
\end{pmatrix}
\]
\marginal{dét. de $\alpha, \beta \mapsto \operatorname{Tr}(\alpha)\operatorname{Tr}(\beta)$}
\note{le « 4 » surcharge un autre chiffre ; la matrice n'est pas complétée.}

\page{94}
Il faudrait prouver que si \uncertain{l'on suppose} qu'il y a une $\Theta \in k.1 \oplus ku \oplus kv$, \uncertain{partout} non nulle, qui est \uncertain{orthogonale} à $1, u, v$ --- il faut des calculs (pour $\operatorname{Tr}(xy)$).
\struck{\uncertain{déterminant}} \uncertain{Le} discriminant de $\operatorname{Tr}$ par rapport à la base $1, u, v$, on trouve
\note{en regard, au centre de la page :}
\[
\text{discriminant de $\operatorname{Tr}$ par rapport à la base $1, u, v$ de $k1 \oplus ku \oplus kv$} = -2\,Q(t;\, b, c, b', c')
\]
[et le discriminant \emph{réduit} de la forme quadratique $\det$ est $-Q(t;\, b, c, b', c')$]
\marginal{Vérifier ça ? --- Oui}

Donc le fait \struck{signifie} \add{implique} que \struck{\ill{}} il y a une $\exists\, \Theta$ orthogonale à $1, u, v$ partout non nulle et \emph{isotrope}.

Donc \struck{\uncertain{pour}}
\[
\det\Theta = \operatorname{Tr}\Theta = 0
\]
Or \struck{l'ensemble des $x$} \struck{\ill{}} les \uncertain{éléments} isotropes dans l'hyperplan $\operatorname{Tr}\Theta = 0$ correspondent aux Borels, \struck{comme \ill{}} on gagne !

\page{95}
Supposons maintenant que
\[
\delta = \delta' = 0 \qquad \text{i.e.} \quad b^2 - 4c = b'^2 - 4c' = 0
\]
ce qui signifie que \struck{\uncertain{$A_1, A_2$ sont isomorphes} \ill{}} \add{$A_1, A_2$} loc. isom. à $k[T]/T^2$. (Je n'ai pas \uncertain{eu le temps} \uncertain{de vérifier dans le cas} \ill{}~: fait évident, \uncertain{quand} \ill{} \ill{} \uncertain{base} générale --- \uncertain{évident} \ill{} \ill{} \uncertain{corps}, plus \uncertain{généralement}~: [si $2$ est inversible \ill{} \ill{}~: si $b \in k$ tel que $b^2 = 0$, $\exists$ loc. $\alpha, \beta \in k$, $\beta \in k^*$, tels que
\[
\alpha^2 = 0, \qquad \beta(2\alpha + b) = 0
\]
\uncertain{Évident} si $2$ \uncertain{inversible}, mais faux si $2.1_k = 0$ (car.~2) car \struck{$\beta b$} dans l'égalité \ill{} $\beta b = 0$ donc (comme $\beta \in k^*$) $b = 0$, mais dans \uncertain{une situation} \uncertain{où} $b^2 = 0$, $b \neq 0$, ça coince~!])
\marginal{\ill{} \ill{} \ill{} \ill{} \uncertain{loc.} \ill{}}

Ça signifie aussi que \struck{Mais dans alg. d'Az.} $A_1, A_2$ \uncertain{sont} des sous-algèbres unipotentes, \struck{\ill{}} \struck{à la position \ill{}} \uncertain{tribu} d'un $\Theta \in M$ (Azumaya \uncertain{ambiante}) qui \uncertain{soit} \struck{de $\operatorname{Tr}$}
\[
\operatorname{Tr}\Theta = \det\Theta = 0
\]
\struck{Soient les \uncertain{unipotents} \ill{} d'un tel} Il y a corr. $1 \to 1$ entre \ill{}
\note{la seconde moitié de la page, à partir de « Mais dans alg. d'Az. », est barrée d'un long trait oblique.}

\page{96}
Donc on \uncertain{introduira} \uncertain{(au cas où)} $A_1 = k.1 + D_1$, $A_2 = k.1 + D_2$, $D_1, D_2$ contenus dans
\[
M^0 = \{x \in M \mid \operatorname{Tr} x = 0\}
\]
et isotropes, i.e.
\[
\operatorname{Tr} u = \operatorname{Tr} v = 0, \qquad \det u = \det v = 0 .
\]
Dans ce cas là, on trouve
\[
Q(t, b, b', c, c') = t^2 .
\]
Mais dans ce cas on \uncertain{là} \uncertain{le} vérifie aussitôt que pour $u \in M$
\[
\operatorname{Tr}(ux) = 0 \iff x \text{ normalise } ku
\]
[i.e. $[x, u] = xu - ux = 0$]
\note{ainsi sur la page : la condition entre crochets est celle de commutation, non de normalisation.}

donc i.e.
\[
x \in B_u \qquad \text{(sous-alg. de Borel normalisatrice de $u$)}
\]
donc si $t = \operatorname{Tr}(uv) = 0$, cela implique que $u, v$ engendrent une sous-alg. de Borel
\[
k \oplus ku \oplus kuv = B_u = B_v
\]
\note{on lit $kuv$ pour le dernier terme ; $kv$ serait attendu.}
Cette \ill{} \ill{}). Donc

\page{97}
\textbf{Proposition}. Si $b = c = b' = c' = 0$, alors
\[
M[b, c, b', c';\, t] = M(0, 0, 0, 0, t) \text{ est Azumaya,}
\]
\struck{ssi $t \in k^*$} ssi
\[
Q(0, 0, 0, 0, t) = t^2 \in k^*
\]
i.e. $t \in k^*$.

En effet, si $t \in k^*$, alors on « réalise » $t$ dans l'algèbre \uncertain{des matrices} $M_2$ p.ex. \uncertain{par} \uncertain{(loc.)}
\[
u = \begin{pmatrix} 0 & 1\\ 0 & 0 \end{pmatrix} \quad \text{et} \quad v = t\begin{pmatrix} 0 & 0\\ 1 & 0 \end{pmatrix}
\qquad
\operatorname{Tr} uv = t \quad \text{OK.}
\]
\ill{} on trouve \ill{} \ill{} \uncertain{Si} $\delta = \delta' = 0$ \ill{}, le \uncertain{cas} précédent \uncertain{se} \ill{} \uncertain{à} \ill{} \ill{} \uncertain{pour} \ill{} \ill{} \ill{} \uncertain{Azumaya}.
\marginal{Il faudrait \ill{} \uncertain{Azumaya}, \ill{} \ill{} \ill{} \ill{} $Q(\ldots) \ldots$ \ill{} \ill{} \uncertain{réduit} \ill{} \ill{}}
\note{une longue note oblique dans la marge gauche, en partie barrée, ne se lit pas au-delà de quelques mots.}

\struck{\uncertain{À}} Reste le cas où $\delta \neq 0$, $\delta' = 0$ \struck{\ill{}}. On peut supposer [un peu court~!!]
\[
b' = c' = 0
\]
donc \uncertain{vérifions}
\[
Q(b, c, b', c', t) = t^2
\]
donc $Q$ inv. \uncertain{signifie} $t$ inv.

\emph{Montrons} que \emph{seul} \uncertain{signifie} que $M(b, c, b', c';\, t)$ est Azumaya.

\page{98}
\textbf{Th.} Soient $b, c, b', c', t \in k$, d'où $M = M(b, c, b', c';\, t)$, avec base can. $\underline{1}, \underline{u}, \underline{v}, \underline{uv}$, et formes $\operatorname{Tr}$ et $\det$ \ldots\ Pour que $M$ soit Azumaya, il faut et suffit que la forme $\det$ sur $k \oplus k\underline{u} \oplus k\underline{v}$ soit non dégénérée, i.e. \uncertain{explicitement} que son discriminant \uncertain{réduit} désigné par $-Q$
\[
\Delta(b, c, b', c', t) = -\bigl(t^2 - bb't + b^2c' + b'^2c - 4cc'\bigr)
\]
soit inv.
\note{le dernier terme est écrit $b'^2_{\,2}$ sur la page ; on lit $b'^2c$ d'après $Q$ (p.~89).}

En termes plus intrinsèques, soit $M \supset A_1, A_2$ \struck{\ill{}}

\uncertain{\struck{\ill{}}} une algèbre \add{\uncertain{assoc.} unit.} \add{loc.} libre de rang 4 sur $k$, \struck{\uncertain{munie de} \ill{}} \add{loc. libres, deux sous-alg. transverses} $A_1, A_2$ \add{de rang 2}, \uncertain{bien} supposons qu'il y ait dans $M$ \struck{\ill{}}

des formes $\operatorname{Tr}$, $\det$ avec les propr. habituelles, \add{(\ill{} \ill{} \ill{} \uncertain{libres}) \ill{}} et que $A_1, A_2$ \uncertain{soient} \struck{\ill{}} \uncertain{étales} sur $k$ (\uncertain{de} \uncertain{rang} \uncertain{loc.}) \ill{} \ill{} $k \cdot 1$ (\uncertain{i.e.} $A_1 + A_2$). \uncertain{\ill{}} \uncertain{Ce} qui implique que $A_1$ \uncertain{engendré} $u$, $A_2$ \uncertain{par} $v$, d'où $b, c, b', c', t$, \ldots\ $M \simeq M(b, c, b', c';\, t)$).

Alors $M$ est azumayenne ssi la restriction de la forme \struck{\uncertain{dét}} \uncertain{quadr.} $\det$ \add{à} $A_1 + A_2$ est non dégénérée,

\page{99}
\emph{Dém.} C'est géométrique, on peut se placer sur un corps alg. clos. Nécessité~: supposons que $M$ soit Azumaya, \add{prouvons non dégénérescence.} \struck{On peut \uncertain{\ill{}}}

Si $A_1, A_2$ sont \uncertain{multiplicatifs}, \ill{} \uncertain{prouvé} le résultat, et \uncertain{aussi} si $A_1, A_2$ sont toutes les deux \uncertain{étales}. Il reste le cas où l'une \uncertain{seule} est \uncertain{étale}, l'autre unipotente, mais il \uncertain{n'y a} qu'une seule position \add{relative}. Dans ce cas on choisit $v$ unipotent, $u$ un projecteur
\[
\text{\struck{$b = 1,\ c = 0$}}, \quad b' = 0, \quad c' = 0
\]
\[
Q = t^2, \qquad \Delta = -Q = -t^2
\]
\ill{} vérifier que $t \neq 0$. On choisit
\note{d'abord, biffé : $v = \begin{pmatrix} 0 & 1\\ 0 & 0 \end{pmatrix}$, $u = \begin{pmatrix} 1 & 0\\ 0 & 0 \end{pmatrix}$ ; non biffé à la suite :}
\[
uv = \begin{pmatrix} 0 & 1\\ 0 & 0 \end{pmatrix}
\]
\[
v = \begin{pmatrix} 0 & 1\\ 0 & 0 \end{pmatrix} \qquad u = \begin{pmatrix} \alpha & \beta\\ \gamma & \delta \end{pmatrix} \quad \text{« général »,}
\]
\[
uv = \begin{pmatrix} 0 & \alpha\\ 0 & \gamma \end{pmatrix} \qquad \operatorname{Tr} uv = \gamma
\]
\[
\delta(u) = (\alpha + \delta)^2 - 4(\alpha\delta - \beta\gamma)
\]
On peut choisir $\alpha, \beta, \gamma, \delta$ tels que $\gamma\,\delta(\alpha, \beta, \gamma, \delta) \neq 0$.

Suffisance~: \struck{Il suffit de} \ldots
\note{la page s'arrête là ; la démonstration de la suffisance n'est pas écrite.}

\note{la démonstration de la suffisance n'est pas poursuivie dans ce lot ; elle sort peut-être du lot.}

\end{document}
